Um die gegebene Fläche als reelle Mannigfaltigkeit vom Geschlecht 2 korrekt zu beschreiben, müssen wir die Verklebung der Ränder präzise definieren und die lokalen Karten sowie die Übergangsfunktionen korrekt formulieren.

Betrachten wir das Rechteck $R = [0, 2] \times [0, 1]$. Die Verklebung der Ränder erfolgt durch die Identifikationen $(x, 0) \sim (x, 1)$ für $x \in [0, 2]$ und $(0, y) \sim (2, 1-y)$ für $y \in [0, 1]$. Diese Verklebung führt zu einer Fläche vom Geschlecht 2.

Wir definieren lokale Karten, die diese Struktur beschreiben:

\begin{align*}
\text{Karte um die Ecke } (0,0): \quad & \varphi_1: U_1 \to \mathbb{R}^2, \quad \varphi_1(x,y) = (x,y), \\
& \text{wobei } U_1 \text{ eine kleine Umgebung von } (0,0) \text{ in } R \text{ ist.} \\
\text{Karte entlang der Kante } [0,2] \times \{0\}: \quad & \varphi_2: U_2 \to \mathbb{R}^2, \quad \varphi_2(x,y) = (x,y), \\
& \text{wobei } U_2 = [0,2] \times (-\epsilon, \epsilon) \text{ für ein kleines } \epsilon > 0. \\
\text{Karte im Inneren des Rechtecks:} \quad & \varphi_3: U_3 \to \mathbb{R}^2, \quad \varphi_3(x,y) = (x,y), \\
& \text{wobei } U_3 \text{ eine offene Menge im Inneren von } R \text{ ist.}
\end{align*}

Die Übergangsfunktionen zwischen diesen Karten sind wie folgt definiert:

\begin{align*}
\varphi_1 \circ \varphi_2^{-1}: \quad & (x,y) \mapsto (x,y), \quad \text{für } (x,y) \in U_1 \cap U_2, \\
\varphi_2 \circ \varphi_3^{-1}: \quad & (x,y) \mapsto (x,y), \quad \text{für } (x,y) \in U_2 \cap U_3, \\
\varphi_3 \circ \varphi_1^{-1}: \quad & (x,y) \mapsto (x,y), \quad \text{für } (x,y) \in U_3 \cap U_1.
\end{align*}

Die Verklebung der Ränder wird durch die Übergangsfunktion $\varphi_2 \circ \varphi_2^{-1}$ beschrieben, die die Identifikation $(x, 0) \sim (x, 1)$ und $(0, y) \sim (2, 1-y)$ berücksichtigt. Diese Übergangsfunktionen sind glatt, da sie Identitäten oder Verschiebungen darstellen.

%CHECK: Korrekte Verklebung der Ränder beschrieben, Übergangsfunktionen angepasst, irrelevante Cauchy-Riemann-Gleichungen entfernt.